\documentclass[../../../include/open-logic-section]{subfiles}

\begin{document}

\olfileid{sth}{replacement}{intro}
\olsection{引言}

替换公理是区分$\ZF$与~$\Z$的公理模式。在
\crefrange{sth:ordinals::chap}{sth:spine::chap}中，我们一直在使用它。在本章中，我们将最终考虑这个问题：替换公理能得到证成吗？

为了使问题更加明确，值得指出的是，替换公理确实相当\emph{强}。我们将在本章（并在\olref[sth][card-arithmetic][fix]{sec}中再次）认识到它究竟有多强。但这恰恰说明，证成确实是必要的。

我们将讨论两种证成。粗略地说：一种\emph{外在的}证成试图依据公理的成效来证成它；一种\emph{内在的}证成则试图表明，公理被所涉及的数学概念本身所证成。我们将在本章中更好地理解这两种证成的含义，但这仅仅是冰山一角。欲了解更多，特别参见\citeauthor{Maddy1988a}（\citeyear{Maddy1988a}与\citeyear{Maddy1988b}）。

\end{document}